%
% Tex input file for CGNS Steering Committee Charter
%
% To generate a DVI file, then a PostScript file named charter.ps:
%
%     latex charter.tex
%     dvips charter.dvi -o
%
% To generate a PDF file named charter.pdf:
%
%     pdflatex charter.tex


\documentclass[twoside,fleqn]{article}

%
% Packages
%
%\usepackage{times}			% use PostScript Times-Roman for text
\usepackage{indentfirst}		% indent first paragraph in sections
%\usepackage[first,light]{draftcopy}	% add DRAFT watermark
\usepackage[tbtags]{amsmath}		% use ams math package
\usepackage{graphicx}			% for including eps files
%\usepackage{subfigure}			% for sub-figures
\usepackage[bf]{caption}		% for more flexible captions
   \setlength{\captionmargin}{18pt}	% caption margins
\usepackage{array}			% use new array and table package
   \setlength{\extrarowheight}{2pt}	% increase row heights in tables
\usepackage{longtable}			% use longtable package
\usepackage{alltt}			% verbatim input with font changes
%\usepackage{lscape}			% for landscape pages
%\usepackage{mdwtab}			% use mdwtab array and table package
%\usepackage{tabls}			% access \hline[dimen] feature
%\usepackage{dcolumn}			% allow table alignment at dec points
%   \newcolumntype{d}[1]{D{.}{.}{#1}}	% for tables, arg is # decimal places
\usepackage{natbib}			% for author-year citations
   \setlength{\bibhang}{0pt}		% don't use hanging indent format
\usepackage[normalem]{ulem}		% allow underlining
\usepackage{calc}			% calculation package
\usepackage{ifthen}			% control structures
\usepackage{fancyhdr}			% fancy headers and footers
   \fancyhf{}
   \renewcommand{\headrulewidth}{0pt}
   \renewcommand{\footrulewidth}{0pt}
\usepackage{xspace}			% common-sense spacing after text macro
\usepackage{titling}			% more title page control
\usepackage{mdwlist}			% more flexible description lists
\usepackage{color}			% colors and grays
%   \definecolor{subcolor}{rgb}{1.0,0.934,0.801}
   \definecolor{subcolor}{rgb}{0.809,0.852,1.0}
%\usepackage[T1]{fontenc}		% for better-looking "_" and tt "{"
\usepackage{ae}				% used instead of above for PDF output
\usepackage{hyperref}			% for hypertext links in PDF
   \hypersetup{letterpaper,plainpages=false,
               pdftitle={The CFD General Notation System Overview and Entry-Level Document},
               pdfauthor={CGNS Project Group},
               colorlinks,
               linkcolor=blue,citecolor=blue,filecolor=red,pagecolor=blue,
               urlcolor=red}
   \renewcommand{\sectionautorefname}{Section}
   \renewcommand{\subsectionautorefname}{Section}
   \renewcommand{\subsubsectionautorefname}{Section}

%
% Page layout
%
\normalsize
\setlength{\oddsidemargin}{0.5in}	% left margins-1.0in for odd/even pages
\setlength{\evensidemargin}{0.0in}
\setlength{\textwidth}{6.0in}		% text width
%\setlength{\marginparwidth}{0in}	% margin note space parameters
%\setlength{\marginparsep}{0in}
%\setlength{\oddsidemargin}{0.25in}	% left margins-1.0in for odd/even pages
%\setlength{\evensidemargin}{0.75in}
%\setlength{\textwidth}{5.5in}		% text width
\setlength{\marginparwidth}{1.1in}	% margin note space parameters
\setlength{\marginparsep}{0.2in}
\reversemarginpar
\setlength{\topmargin}{-0.5in}		% top margin-1.0in
\setlength{\headheight}{0.25in}		% header space parameters
\setlength{\headsep}{0.5in}
\setlength{\textheight}{8.5in}		% text height
\setlength{\topskip}{\baselineskip}	% dist from top of body to 1st baseline
\setlength{\footskip}{0.75in}		% dist from bottom of body to footer
\raggedbottom				% to avoid stretching vertical space
\pagestyle{fancy}			% using fancyhdr package
\fancypagestyle{plain}{%		% page # in corner for plain style
   \fancyhf{}
   \fancyfoot[LE,RO]{\bfseries \thepage}}

%
% Misc spacing parameters
%
%\setlength{\parskip}{\baselineskip}	% blank space between paragraphs
\setlength{\parskip}{0.5\baselineskip}	% blank space between paragraphs
\setlength{\doublerulesep}{0pt}		% for wider \hlines
\setlength{\fboxsep}{5pt}		% for more margin in boxes
\newlength{\saveparindent}		% save for use where LaTeX changes them
\newlength{\saveparskip}		% save for use where LaTeX changes them
\newlength{\savebaselineskip}
\setlength{\saveparindent}{\parindent}
\setlength{\saveparskip}{\parskip}
\setlength{\savebaselineskip}{\baselineskip}
\newlength{\tmplength}			% for temp use wherever needed
\newlength{\Pwidth}			% for use in tables
\newlength{\Pwidtha}

%
% User-defined environments
%
% "Definition" list with variable-length entries. Requires
% mdwlist package. (See LaTeX Companion, p 64, and mdwlist documentation.)
\newenvironment{Ventry}[1]%
   {\settowidth{\tmplength}{#1\hspace{1em}}%
    \begin{basedescript}{%
           \desclabelwidth{\tmplength}%
           \desclabelstyle{\nextlinelabel}%
           \renewcommand{\makelabel}[1]{##1\hfil}}}
   {\end{basedescript}}
% "Definition" list with variable-length entries, indented. Requires
% mdwlist package. (See LaTeX Companion, p 64, and mdwlist documentation.)
\newenvironment{Ventryi}[1]%
   {\settowidth{\tmplength}{\hspace{\saveparindent}#1\hspace{1em}}%
    \begin{basedescript}{%
           \desclabelwidth{\tmplength}%
           \desclabelstyle{\nextlinelabel}%
           \renewcommand{\makelabel}[1]{\hspace{\saveparindent}##1\hfil}}}
   {\end{basedescript}}
%
% Adapted from comp.text.tex post by Keith Reckdahl
% (reckdahl@am-sparc7.Stanford.EDU)
% Indents text from the left margin by current value of the list
% length \leftmargin
\newenvironment{indleft}%
   {\begin{list}{}
      {\setlength{\topsep}{0pt}%
       \setlength{\listparindent}{0pt}%
       \setlength{\itemindent}{0pt}%
       \setlength{\parsep}{\parskip}%
      }%
      \item[]}%
   {\end{list}}
% As above, but surrounding an alltt environment
\newenvironment{indlefttt}%
   {\begin{list}{}
      {\setlength{\topsep}{0pt}%
       \setlength{\listparindent}{0pt}%
       \setlength{\itemindent}{0pt}%
       \setlength{\parsep}{\parskip}%
      }%
      \item[]%
      \begin{alltt}}%
   {\end{alltt}\end{list}}
% As above, but also with a tabbing environment
\newenvironment{tabindlefttt}%
   {\begin{list}{}
      {\setlength{\topsep}{0pt}%
       \setlength{\listparindent}{0pt}%
       \setlength{\itemindent}{0pt}%
       \setlength{\parsep}{\parskip}%
      }%
      \item[]%
      \begin{alltt}\begin{tabbing}}%
   {\end{tabbing}\end{alltt}\end{list}}

% "changes" environment, to identify changed code by color.  The color is
% given by the argument, with a default of red.  When the environment is
% exited, the color is changed to black.
\newenvironment{changes}[1][red]%
   {\color{#1}}
   {\color{black}}

%
% User-defined commands
%

% Degree symbol
\newcommand{\degree}[1][]{\ensuremath{{#1}^\circ}\xspace}

% Degree Rankine symbol
\newcommand{\degR}{\textsuperscript{$\circ$}\textnormal{R}\xspace}

% Re-def longtable's caption command to use \captionlabelfont from caption pkg
% (basis lifted from longtable.sty)
\makeatletter
\renewcommand{\LT@makecaption}[3]{%
  \LT@mcol\LT@cols c{\hbox to\z@{\hss\parbox[t]\LTcapwidth{%
    \sbox\@tempboxa{{\captionlabelfont #1{#2: }}#3}%
    \ifdim\wd\@tempboxa>\hsize
      {\captionlabelfont #1{#2: }}#3%
    \else
      \hbox to\hsize{\hfil\box\@tempboxa\hfil}%
    \fi
    \endgraf\vskip\baselineskip}%
  \hss}}}
\makeatother

% Make next page odd, with preceding blank page empty (LaTeX Companion, p 93)
\newcommand{\clearemptydoublepage}{\newpage{\pagestyle{empty}\cleardoublepage}}

% Text subscripts, analogous to \textsuperscript command (from comp.text.tex
% post by rf@cl.cam.ac.uk (Robin Fairbairns))
\makeatletter
\DeclareRobustCommand*\tsub[1]{%
  \@tsub{\selectfont#1}}
\def\@tsub#1{%
  {\m@th\ensuremath{_{\mbox{\fontsize\sf@size\z@#1}}}}}
\makeatother

% Text superscripts, but shorter
\newcommand{\tsup}[1]{\textsuperscript{#1}}

% Shortcuts for specific fonts
\newcommand{\bold}[1]{{\normalfont\textbf{#1}}}  % Bold
\newcommand{\ital}[1]{{\normalfont\textit{#1}}}  % Italic
\newcommand{\key} [1]{{\normalfont\texttt{#1}}}  % Fixed font for keywords
\newcommand{\fort}[1]{{\normalfont\texttt{#1}}}  % Fixed font for Fortran stuff

% Shortcuts for math character sizes
\newcommand{\Dstyle}{\displaystyle}
\newcommand{\Tstyle}{\textstyle}
\newcommand{\Sstyle}{\scriptstyle}

% Use section titles as marks (i.e., in headers) instead of subsection titles
\renewcommand{\sectionmark}[1]{\markright{\thesection\ \ #1}}
\renewcommand{\subsectionmark}[1]{}

% "Better" treatment of headings for Appendices  (LaTeX Companion, pp 29,30)
%\renewcommand{\appendix}{%
%   \newcommand{\app}{%
%      \secdef\Appendix\sAppendix}%
%   \newcommand{\subapp}{%
%      \secdef\subAppendix\subsAppendix}%
%   \newcounter{app}%
%   \newcounter{subapp}[app]%
%   \renewcommand{\theapp}{\Alph{app}}%
%   \renewcommand{\thesubapp}{\theapp.\arabic{subapp}}%
%   \renewcommand{\sectionmark}[1]{\markright{\theapp\ \ ##1}}
% }
%  Use just "Appendix.", without the letter "A" 
%\newcommand{\Appendix}[2][?]{%   Complex form
%   \refstepcounter{app}%
%   \addcontentsline{toc}{section}%
%      {\protect\bfseries{\appendixname~\theapp.\ \ } #1}%
%   {\noindent\Large\bfseries\appendixname\ \theapp.\ \ #2\par}%
%   \sectionmark{#1}\vspace{\baselineskip}}
%\newcommand{\Appendix}[2][?]{%   Complex form
%   \refstepcounter{app}%
%   \addcontentsline{toc}{section}%
%      {\protect\bfseries{\appendixname.\ \ } #1}%
%   {\noindent\Large\bfseries\appendixname.\ \ #2\par}%
%   \sectionmark{#1}\vspace{\baselineskip}}
%\newcommand{\sAppendix}[1]{%   Starred form
%   \refstepcounter{app}%
%   {\noindent\Large\bfseries\appendixname.\ \ #1\par}%
%   \vspace{\baselineskip}}
%
%\newcommand{\subAppendix}[2][?]{%   Complex form
%   \refstepcounter{subapp}%
%   \addcontentsline{toc}{subsection}%
%       {\thesubapp\ \ #1}%
%   \vspace{\baselineskip}
%   {\noindent\large\bfseries\thesubapp\ \ #2\par}%
%   \vspace{\baselineskip}}
%  Use just the subappendix name, without the prefix "A.1", etc.
%\newcommand{\subsAppendix}[1]{%   Starred form
%   \refstepcounter{subapp}%
%   \vspace{\baselineskip}
%   {\noindent\large\bfseries\thesubapp\ \ #1\par}%
%   \vspace{\baselineskip}}
%\newcommand{\subsAppendix}[1]{%   Starred form
%   \refstepcounter{subapp}%
%   \vspace{\baselineskip}
%   {\noindent\large\bfseries#1\par}%
%   \vspace{\baselineskip}}

\begin{document}

\pagenumbering{roman}

%\fancyhead[LE]{\bfseries CGNS Steering Committee Charter}
%\fancyhead[RO]{\bfseries \rightmark}
\fancyhead[LE,RO]{\bfseries CGNS Steering Committee Charter}
\fancyfoot[LE,RO]{\bfseries \thepage}

%\setlength{\droptitle}{1.0in}
\pretitle{\begin{flushleft}\LARGE%
          \includegraphics[width=2.0in]{../logo/cgns_bw}\\*[0.25in]}
\posttitle{\par\end{flushleft}\vskip 1.0em}
\title{{\bfseries CFD General Notation System\\
Steering Committee Charter}\\*[0.25in]
{\Large Document Version 3.1.2}}
\author{}
\date{}
\maketitle
\thispagestyle{empty}

%\title{{\LARGE \textbf{CGNS}}\\*[0.5in]
%The CFD General Notation System\\
%Steering Committee Charter\footnote
%{This document was originally written by
%Bob Bush (Pratt \& Whitney),
%Raymond Cosner (Boeing, Phantom Works),
%Diane Poirier (ICEM CFD Engineering),
%Ed Hall (Rolls Royce/Allison),
%Armen Darian (Boeing, Rocketdyne), and
%Dan Dominik (Boeing, Reusable Space Systems)
%for the CGNS Project Group.}}
%\author{CGNS Project Group}
%\date{}
%\maketitle
%\thispagestyle{empty}

\clearemptydoublepage
\setlength{\parskip}{0ex}		% remove blank space between paragraphs
\thispagestyle{plain}
\tableofcontents
\setlength{\parskip}{\saveparskip}	% put it back

\clearemptydoublepage
\pagenumbering{arabic}
\section{CGNS Background}
\label{s:background}
\thispagestyle{plain}

The information in this section is supplied for the sake of
completeness. It is identical with the information found in the
Overview.

\subsection{Purpose}

The purpose of CGNS is to provide a standard for recording and
recovering computer data associated with the numerical solution of the
equations of fluid dynamics. The format implemented by this standard is
(1) general, (2) portable, (3) expandable, and (4) durable.

The CGNS system consists of a collection of conventions, and software
implementing those conventions, for the storage and retrieval of
CFD (computational fluid dynamics) data. The system consists of two
parts: (1) a standard format for recording the data, and (2) software
that reads, writes and modifies data in that format.  The format is a
conceptual entity established by the documentation; the software is a
physical product supplied to enable developers to access and produce
data recorded in that format. The CGNS standard, applied through the use
of the supplied software, is intended to do the following:

\begin{itemize*}
\item facilitate the exchange of CFD data
      \begin{itemize*}
      \item between sites.
      \item between applications codes.
      \item across computing platforms.
      \end{itemize*}
\item stabilize the archiving of CFD data.
\end{itemize*}

\subsection{Participation and Brief History}

The CGNS project originated around 1994--1995 through a series of
meetings between Boeing and NASA that addressed improved means for
transferring NASA technology to industrial use. It was held that
a principal impediment to technology transfer was the disparity
in I/O formats employed by various flow codes, grid generators,
and so forth. The CGNS system was conceived as a means to promote
``plug-and-play'' CFD.

Agreement was reached to develop CGNS at Boeing, under NASA Contract
NAS1-20267, with active participation by a team of CFD researchers from

\begin{itemize*}
\item NASA Langley Research Center
\item NASA Glenn Research Center
\item NASA Ames Research Center
\item Boeing St-Louis (McDonnell-Douglas Corporation).
\item Boeing Commercial Airplane Group Aerodynamics
\item Boeing Commercial Airplane Group Propulsion
\item ICEM CFD Engineering Corporation of Berkeley, California
\end{itemize*}
Also participating in the discussions at various times have been
researchers from
\begin{itemize*}
\item Defense and Space Group and Environmental Systems
\item Arnold Engineering Development Center, representing the NPARC Alliance
\item Wright-Patterson Air Force Base
\end{itemize*}

\subsection{Scope}

The principal target of CGNS is the data normally associated with
compressible viscous flow (i.e., the Navier-Stokes equations), but the
standard is also applicable to subclasses such as Euler and potential
flows.

CGNS Version 1.0, released 5/15/98, was limited to problems described
by multiblock structured grids. Version 1.1 addresses grids,
flowfields, boundary conditions, and block-to-block connection
information. Also included are a number of auxiliary items,
including nondimensionalization, reference state, and equation set
specifications. The extension to time-dependent flows and unstructured
grids is addressed in Version 2. Also included are links between CGNS
data and CAD geometry. Any mix of the following types of field data can
be recorded:

\begin{itemize*}
\item nodal.
\item edge-centered.
\item face-centered.
\item cell-centered.
\end{itemize*}
Block connections can be of the following types:
\begin{itemize*}
\item contiguous (one-to-one).
\item abutting (patched mismatched).
\item overlapping (Chimera).
\end{itemize*}

Much of the standard and the software is applicable to computational
field physics in general. Disciplines other than fluid dynamics would
need to augment the data definitions and storage conventions, but the
fundamental database software, which provides platform independence, is
not specific to fluid dynamics.


\clearemptydoublepage
\section{Mission/Vision/Responsibilities}
\label{s:mission}
\thispagestyle{plain}

The mission of the CGNS Steering Committee is to ensure the continuation
of the CFD General Notation System.
The survival of a standard depends entirely on its level of use.
Therefore the CGNS Steering Committee must aim at providing a standard
that is widely accepted by the CFD community.

Several elements must be satisfied to ensure the acceptance of the
CGNS standard.
The most obvious asset is that CGNS must be useful.
Not only must it answer the current needs for the recording of fluid
dynamics data, but it must also follow the changes in requirements as
CFD progresses.
A second important element is that CGNS must be easy to implement.
The CGNS Mid-Level Library (or
Application Programming Interface, API) must be user-friendly and
well documented, and online
support must be available for all users at all times.
The standard must also be easily accessible, meaning that all the
sources, binaries, documents and any other pertinent information must be
available to anyone without restrictions.
Finally, it is of utmost importance that CGNS retains its public nature,
encouraging contributions from all users.

The Steering Committee has the responsibility to oversee that the
CGNS standard remains useful, accessible, easy to use, and preserves
its public nature.
This implies multiple activities, which can be subdivided in the
following groups:

\subsection{Ensure the maintenance of the existing software, documentation and web site}
\label{s:maintenance}

The CGNS Steering Committee is responsible for appointing a prime
source, and overseeing the prime source activities.
The Steering Committee must ensure that the prime source maintains the
existing software, documentation and web site.
This includes, but is not limited to:

\begin{itemize}
\item correcting/updating the documentation if necessary
\item fixing any reported software bug
\item collecting a list of CGNS users via the web site
\item keeping the web site up to date with the latest versions of the
      documentation and
      software
\item informing the user base of new releases and major software problems
\item posting proposals for
      new features or modifications to the CGNS standard on the web site
      and collect comments from the user base
\item maintaining a distribution site for
      contributed software
      utilities which utilize the CGNS standard
\end{itemize}

\subsection{Provide mechanisms for the evolution of the standard}
\label{s:evolution}

The CGNS Steering Committee has the responsibility to support and even
encourage the evolution of the standard in order for CGNS to remain
useful.
Therefore, the committee must solicit technical support and ``in-kind''
contributions.
In addition, the Steering Committee must follow the policies described
in \autoref{s:changes} of this document
regarding the collection and evaluation of technical
proposals.

\subsection{Promote the acceptance of CGNS}
\label{s:acceptance}

The CGNS Steering Committee has the responsibility to promote the
acceptance of CGNS throughout the CFD community.
This can be achieved through various means, including word of mouth,
advertising, business articles, and presentations at conferences and
technical meetings.

\subsection{Provide mechanism for answering questions and exchanging ideas}
\label{s:services}

Electronic mail constitutes the main point of contact between CGNS
users and CGNS developers.
Therefore, the CGNS Steering Committee must maintain an electronic mail
forum, to which users can post questions, answer questions, and exchange
ideas.
Members of the CGNS Steering Committee and/or appointed qualified
persons will respond to the posted questions on the forum.

\subsection{Determine the means by which the CGNS activities are supported}
\label{s:support}

The CGNS Steering Committee has the obligation to determine the means by
which all CGNS activities are supported.
The Committee is also responsible for identifying and obtaining sources
of funding, if appropriate.
Finally the CGNS Steering Committee has the responsibility to distribute
the tasks and funds to the most appropriate candidate, in the best
interests of CGNS.


\clearemptydoublepage
\section{Organization/Bylaws}
\label{s:organization}
\thispagestyle{plain}

The CGNS Steering Committee is a voluntary organization that will
determine its own policies and internal structure, and will govern by
consensus whenever possible.
In the absence of consensus, a two-thirds majority of the Committee
members will be required to adopt changes to the standard, alter this
Charter, or take other official actions.

The CGNS Steering Committee will meet at a minimum of one time per year.
The time and location will be determined by consensus of the Committee,
and all members of the Committee will be notified in advance.

The members of the CGNS Steering Committee will appoint a Chairperson
whose responsibilities will include coordinating activities,
facilitating meetings and serving as a focal point for the Committee.
The Chairperson will be a member of the Committee, be elected by
consensus, and serve for a two-year term.
There is no limit on how many terms the Chairperson can be elected.
At the discretion of the Chairperson, a Vice-Chairperson may be
appointed by consensus of the Committee, to assist the Chairperson with
his or her duties.
The Vice-Chairperson will also be a member of the Committee.
The appointment of a secretary to maintain records will be at the
discretion of the Chairperson.

The CGNS Steering Committee may decide to suggest appropriate
contributions from its members.
The Steering Committee is not prohibited from charging membership fees;
the decision whether to do so, and the amount of the fees, lies within
the purview of the Steering Committee.

All parties are welcome to bring forward issues and participate in
development of the CGNS Standard, whether or not they are members of the
Steering Committee.

The decision whether to support the migration of the CGNS standard to
ISO/STEP, or any other organization, lies within the purview of the
Steering Committee.

\subsection{Representation}
\label{s:representation}

The CGNS Steering Committee will be made up of representatives from
specific institutions, rather than individuals.
Changes or additions to Steering Committee membership will be based on
potential contribution to the standard.
Membership on the Steering Committee will be limited to 30 institutions
that actively participate in the development, maintenance, distribution
and use of the CGNS Standard.
No more than 5 institutions may be related, i.e., have the same parent
organization.
Changes to the Membership (including the limit on the number of
institutions) will be determined by consensus, or if required, a
two-thirds majority of the existing Membership.

To help satisfy the duties of the Steering Committee as a whole, as
described in \autoref{s:mission}, the minimal responsibilities of each
individual Steering Committee member are to:
\begin{itemize*}
\item Attend as many telecons/meetings as possible, but not less than
      one per year
\item Read and send comments on proposals or other issues when asked to
      do so
\item Vote when asked to do so
\end{itemize*}
More active participation --- including support, software development,
and actively working to improve and promote CGNS --- is encouraged.

\newpage
The Steering Committee members as of 14 April 2011 are:

\begin{itemize*}
\item ADAPCO
\item ANSYS-CFX
\item ANSYS-ICEM CFD
\item Aerospatiale Matra -- Airbus
\item Boeing Commercial
\item Computational Engineering Solutions
\item Concepts NREC
\item Intelligent Light
\item NASA Glenn
\item NASA Langley
\item ONERA
\item Pointwise, Inc.
\item Pratt \& Whitney
\item Rolls-Royce / Allison
\item Stanford University
\item Stony Brook University
\item Tecplot, Inc.
\item TTC Technologies
\item University of Colorado
\item U. S. Air Force / AEDC
\end{itemize*}

\subsection{Standing Committees}
\label{s:committees}

The CGNS Steering Committee may constitute Standing Committees, in
an ongoing or temporary basis, to which it may delegate various
responsibilities.
The Standing Committees will report and make recommendations to the
Steering Committee who will retain the authority to act and make final
decisions.

\subsection{Software and Documentation Support Team}
\label{s:supportteam}

The CGNS Steering Committee will be responsible for selecting one or
more organizations to maintain and distribute existing documentation
and software, to develop and distribute new software resulting from
extensions to the standard, and to post or distribute meeting minutes
and other new documentation.

The organization(s) selected to maintain CGNS software will determine
the form of newly developed software and maintain compatibility with the
existing ADF Core and CGNS API.

The organization(s) selected to maintain CGNS Documentation will
be responsible for posting and maintaining on the web the Steering
Committee meeting minutes, Charter, and archive information.


\clearemptydoublepage
\section{Standard and Software Governing Principles}
\label{s:principles}
\thispagestyle{plain}

\subsection{Distribution}
\label{s:distribution}

This section describes the policy governing the distribution of the
CGNS standard and software to the engineering and scientific community
at large.
By definition, the CGNS \textit{standard} refers to the
Standard Interface Data Structures
(SIDS) definitions, the SIDS File Mapping, and the CGNS Mid-Level
Library structure (API), as well as all documentation.
The CGNS \textit{software} refers to the CGNS
Library source code, the CGIO core routines,
and the ADF and HDF5 database manager implementations.
The CGNS \textit{software} may also include sample programs
demonstrating the application and use of the CGNS and ADF libraries,
as well as some utility programs to assist with the implementation and
analysis of CGNS-based files and systems.

Implementation and maintenance of the CGNS distribution policy is the
responsibility of the CGNS Steering Committee.
The distribution policy dictates that both the CGNS standard and the
CGNS software are publicly available, and that the standard and software
itself are free of charge.
The CGNS software may be used for any purpose, including commercial
applications, and may be altered and redistributed, subject to the
restrictions described in the CGNS License (see the
\hyperref[s:license]{Appendix}).

It is the responsibility of the CGNS Steering Committee to enable
distribution mechanisms that comply with the following principles:

\begin{itemize}
\item The CGNS standard (documentation and definitions) will be publicly
      available at no more than the cost of distribution.
\item The CGNS software (CGNS Library and CGIO, ADF and HDF5 core), including source
      code, will also be available at no more than the cost of
      distribution.
\item The CGNS API (Mid-Level Library), including source code will be
      similarly available.
\item Development, sale, and licensing of proprietary packages based on
      CGNS that perform substantive operations on the data, beyond the
      I/O performed by the API, are encouraged.
      Such packages must abide by the restrictions described in the
      CGNS License.
\item The sale of services designed to assist in the conversion of
      existing software to the CGNS standard is acceptable.
\item The voluntary contribution of software that performs operations on
      CGNS data is encouraged.
\item The CGNS Steering Committee will provide mechanisms for the
      accumulation and distribution of contributed software, but will not
      be responsible for the function of contributed software.
\item Contributed software does not become part of the CGNS Standard,
      that is, either the SIDS or the API, without the approval of the
      CGNS Steering Committee.
\item The Steering Committee may agree to support or endorse additional
      utility software.
\item The Steering Committee will \emph{not} endorse third party
      software.
\end{itemize}

\subsection{Changes or Additions to the Standard}
\label{s:changes}

CGNS is a standard that has been developed with the key concepts of
flexibility and extendibility in mind.
The standard can accommodate the majority of CFD data quantities in
practical usage today; however, some additional capabilities are still
being implemented.
It is also understood that in the future other additional capabilities
will need to be implemented as well.
For these reasons, a process for adding to or modifying the existing
CGNS standard is necessary.

To address a particular need or deficiency in CGNS, a proposal for a
potential change to the standard first must be made.
A Technical Team will prepare all proposals.
A Technical Team may voluntarily submit the proposal, or a Technical
Team may be specifically appointed by the Steering Committee to author
the proposal.

A primary requirement of all proposals for modifications will be to
support and maintain code compatibility.
No additions or changes to the CGNS standard will be adopted which
invalidate existing software or data.

Prior to adoption, the Technical Team must present all proposals in
an open and public forum.
Included with the proposal, a draft of the necessary changes to the
SIDS and File Mapping must be provided by the team introducing the
modifications.
The open forum will then review the proposal, identify any possible
shortcomings, and suggest alternatives or improvements.

After the proposal has been presented and deliberated upon, only the
Steering Committee has final authority of approval and may elect to
do one of three things.
First, the Steering Committee may vote by consensus (or a two-thirds
majority if necessary) to accept the proposal as is, and thus the
changes are approved for implementation.
If such approval does not occur, the Steering Committee may still feel
there is merit to the proposal, and may choose to defer acceptance of
the proposal under the provision that specific changes be made.
Finally, the Steering Committee may deem there is little merit in the
proposed changes to CGNS, and reserves the right to reject the proposal
outright.
Whatever the disposition of the proposal, individual organizations may
implement \texttt{UserDefined} functions, provided that they adhere to
the conventions and standards as defined in the SIDS.


\clearemptydoublepage
\appendix
%\fancyhead[RO]{\bfseries \rightmark}
%\fancyfoot[LE,RO]{\bfseries \thepage}
%\app{CGNS License}
\section{CGNS License}
\label{s:license}
\thispagestyle{plain}

The distribution and use of the CGNS software is covered by the
following license:

\begin{indleft}
This software is provided ``as-is'', without any express or implied
warranty.
In no event will the authors be held liable for any damages arising from
the use of this software.

Permission is granted to anyone to use this software for any purpose,
including commercial applications, and to alter it and redistribute it
freely, subject to the following restrictions:
\begin{enumerate}
\item The origin of this software must not be misrepresented; you must
      not claim that you wrote the original software.
      If you use this software in a product, an acknowledgment in the
      product documentation would be appreciated but is not required.
\item Altered source versions must be plainly marked as such, and must
      not be misrepresented as being the original software.
\item This notice may not be removed or altered from any source
      distribution.
\end{enumerate}
\end{indleft}

This license is borrowed from the zlib/libpng License (see
\url{http://www.opensource.org/licenses/zlib-license.php}), and
supercedes the GNU Lesser General Public License (LGPL) which previously
governed the use and distribution of the software.

For details on the policy governing the distribution of the CGNS
standard and software see \autoref{s:principles}.


\end{document}
